\section{Limits of the method}
\label{sec:limits}

\subsection{One step above the counting bound}

Passing from \(k\) to \(k+1\) increases \(\Lambda_0\) by \(k(q-1)+1\),
and by Lemma~\ref{lem:admissible-shape} each run of the admissible set
then has about \(\sqrt{2k}\) elements, while the gap between the roots
of \(Q\) shrinks.  In every scanned case the two runs merge at the first
step, and once \(\cS\) is an interval containing the means \(k\theta_{d-1}/\theta_d\)
and \(R/t\), both integer systems are feasible and every bracket bound is
zero.  For \(\PG(4,7)\) and \(k=32\), for example, \(\Lambda_0=207\) and
\(Q(s)=\tfrac12(7s^2-67s+192)\) has negative discriminant, so \(Q>0\)
everywhere, and \(\cS=\{0,1,\dots,12\}\).  The mechanism therefore gains
one step and no more.  A larger gain would need a constraint that
survives an interval of admissible sizes, and the natural candidate is
the distribution of the losses \(Q(s_\Pi)\) themselves, which the
character equations do not see.

\subsection{Complete caps without coplanar quadruples}

If \(c_4(A)=0\), every point off \(A\) lies on at most one secant.  For
a complete cap this forces index one everywhere, \(\Lambda_0=0\), and
\(Q(s_\Pi)=0\) for every hyperplane: the section sizes are the integer
roots of \(Q\).  In the coding dictionary \(C_A^\perp\) is then a perfect
two-error-correcting code, and Pavese's classification of the equality
cases of his counting bound for \(4\)-general
sets~\cite[Proposition~3.1]{Pavese2025} identifies \(A\) as the elliptic
quadric of \(\PG(3,2)\) or the \(11\)-cap of \(\PG(4,3)\).  The identity
accommodates both: for \(k=5\) in \(\PG(3,2)\), \(Q(s)=(s-1)(s-3)\), and
every plane meets the quadric in \(1\) or \(3\) points; for \(k=11\) in
\(\PG(4,3)\), \(Q(s)=\tfrac32(s-2)(s-5)\), and every solid meets the cap
in \(2\) or \(5\) points, the weights \(9\) and \(6\) of the dual ternary
Golay code.  Complete \(4\)-general sets that are not complete caps exist
in many other parameters~\cite[Section~5]{Pavese2025}, and normal
rational curves have \(c_4=0\) in every dimension; the constraints of
this paper say nothing about them because they are not complete caps.

\subsection{Why degree constraints alone cannot force spread collisions}

The higher-dimensional problem of~\cite{RuddSecantDefects} asks whether a
positive proportion of the secants of a complete cap in \(\PG(d,q)\),
\(d\ge4\), must satisfy \(T_\ell\ge cq\).  The constraints available
without an ambient hyperplane count are the incidence identities of
Section~\ref{sec:preliminaries}: the secant--point
count~\eqref{eq:first-moment}, the collision
count~\eqref{eq:secant-collisions}, the pencil identities~\eqref{eq:pencil},
the coincidence count \(\sum_{x}\binom{r(x)}{2}=3c_4(A)\), and the
capacity \(r(x)\le m\).  We show that formal data can satisfy all of
them with almost every secant free.

\begin{proposition}[Degree constraints do not force spread collisions]
\label{prop:degree-countermodel}
Fix \(d\ge4\) and let \(q\to\infty\) with \(k=k(q)\) satisfying
\(k\asymp q^{(d-1)/2}\), \(\Lambda_0\ge0\), \(\Lambda_0\equiv0\pmod3\),
and \(\Lambda_0\le k(k-1)/24\).  Then there exist an abstract \(k\)-set
\(A\), a family \(\mathcal Q\) of \(M=\Lambda_0/3\) four-subsets of \(A\)
called coplanar, a finite set \(Y\) with \(|Y|=\theta_d-k\), and for each
pair \(\ell\subseteq A\) a \((q-1)\)-subset \(Y_\ell\subseteq Y\), such
that, with \(r(x)=\#\{\ell:x\in Y_\ell\}\),
\(T_\ell=\sum_{x\in Y_\ell}(r(x)-1)\), and \(c_4=|\mathcal Q|\):
\begin{enumerate}[label=(\alph*),nosep]
\item \(1\le r(x)\le2\le m\) for all \(x\in Y\), \(\sum_xr(x)=N(q-1)\),
 and \(\sum_x\binom{r(x)}{2}=3c_4\);
\item two pairs with disjoint endpoints share a point of \(Y\) if and
 only if their union is in \(\mathcal Q\), and then they share exactly one
 point; pairs with a common endpoint share no point;
\item \(T_\ell=\#\{B\in\mathcal Q:\ell\subseteq B\}\) and
 \(\sum_\ell T_\ell=6c_4\), and for every \(\ell\) there are nonnegative
 integers \(z_{\ell,1},\dots,z_{\ell,\theta_{d-2}}\) with
 \(\sum_iz_{\ell,i}=k-2\) and \(\sum_i\binom{z_{\ell,i}}{2}=T_\ell\);
\item \(T_\ell\in\{0,1\}\) for every \(\ell\), so at most
 \(6M=2\Lambda_0=o(N)\) pairs have \(T_\ell>0\).
\end{enumerate}
The hypotheses hold for all large \(q\) whenever \(k\) is within a
bounded distance of the smallest integer satisfying the counting bound
and \(\Lambda_0\equiv0\pmod3\).
\end{proposition}

\begin{proof}
A maximal family of edge-disjoint four-subsets of \(A\), regarded as
copies of \(K_4\) in \(K_k\), has at least \(k(k-1)/72\) members: every
four-subset shares an edge with some member, and a fixed member shares an
edge with at most \(6\binom{k-2}{2}\) four-subsets, itself included, so
the family has at least \(\binom{k}{4}/6\binom{k-2}{2}=k(k-1)/72\)
members.  Choose \(M=\Lambda_0/3\le k(k-1)/72\) of them as \(\mathcal Q\).
For each \(B=\{a,b,c,d\}\in\mathcal Q\) introduce three new points
\(x_{ab|cd},x_{ac|bd},x_{ad|bc}\) and put each into the sets of its two
pairs; since the members of \(\mathcal Q\) are edge-disjoint, each pair
lies in at most one member, so each pair receives at most one such point.
Complete every \(Y_\ell\) with private points to size \(q-1\).  The
shared points have \(r=2\) and there are \(3M\) of them; the private
points have \(r=1\) and there are \(N(q-1)-6M\) of them; so
\(|Y|=N(q-1)-3M=N(q-1)-\Lambda_0=\theta_d-k\), and (a) and (b) hold by
construction.  A pair in a member of \(\mathcal Q\) has \(T_\ell=1\), all
others \(T_\ell=0\), which gives (d) and the first half of (c).  For the
pencil rows take one entry \(2\) (for the pair \(\{c,d\}\)) and \(k-4\)
entries \(1\) when \(T_\ell=1\), and \(k-2\) entries \(1\) when
\(T_\ell=0\), padded with zeros; this needs \(\theta_{d-2}\ge k-2\), which
holds for large \(q\) because \(\theta_{d-2}\ge q^{d-2}\) and
\(k\asymp q^{(d-1)/2}\) with \(d\ge4\).  Finally, near the counting bound
\(\Lambda_0=O(kq)\), and \(kq=o(k^2)\) since \(k\asymp q^{(d-1)/2}\) with
\(d\ge4\); so \(\Lambda_0\le k(k-1)/24\) for large \(q\).
\end{proof}

The data of the proposition are not a cap: no ambient space is present,
and the two-point sets \(Y_\ell\) are not lines.  Its role is to show
that any argument forcing a positive proportion of secants to carry
\(T_\ell\ge cq\) must use more than the identities (a)--(c).  The
hyperplane loss identity is such an additional constraint: it ties the
distribution of the indices to the hyperplane sections, which the formal
data do not have.  Whether it can be pushed beyond the one-step gain of
Section~\ref{sec:exclusions} is the question we leave open.
